% paper_10_nuclear_lattice.tex — GeoVac Paper 10
% The Nuclear Lattice: Discrete Graph Structures for Molecular
% Vibration and Rotation
% Status: Draft, March 12, 2026
% Target: arXiv (physics.comp-ph / quant-ph)
\documentclass[aps,pra,twocolumn,superscriptaddress,longbibliography]{revtex4-2}

\usepackage{amsmath,amssymb,amsthm,graphicx,xcolor,bm,braket,dcolumn,booktabs}
\usepackage{hyperref}

\newtheorem{theorem}{Theorem}
\newtheorem{corollary}{Corollary}
\newtheorem{proposition}{Proposition}

\begin{document}

\title{The Nuclear Lattice: Discrete Graph Structures for Molecular\\
Vibration and Rotation}
\thanks{Paper 10 in the Geometric Vacuum series}

\author{J.~Loutey}
\affiliation{Independent Researcher, Kent, Washington}

\date{March 12, 2026}

%% ========================================================================
%% ABSTRACT
%% ========================================================================

\begin{abstract}
The GeoVac framework derives electronic structure from discrete graph
topology on the unit three-sphere $S^3$: electrons occupy a
momentum-space lattice whose graph Laplacian eigenvalues reproduce the
Schr\"odinger spectrum without continuous spatial integration.  We
extend this program to nuclear degrees of freedom.  Whereas the
electron lattice lives in momentum space (with position emerging as the
dual), the nuclear lattice lives in position space (with momentum
emerging as the dual).  The rotational paraboloid---states
$|J, M_J\rangle$ with $0 \leq J \leq J_{\max}$ and
$-J \leq M_J \leq J$---has \emph{identical} algebraic structure to the
electron $(l, m)$ lattice, both being SO(3) angular momentum
representations with $2J+1$ degeneracy per shell and $\pi$-spacing in
the azimuthal direction.  The vibrational chain exploits the SU(2)
algebraic structure of the Morse oscillator: states
$v = 0, 1, \ldots, v_{\max}$ map to angular momentum projections
$m = j - v$ where $j = \omega_e/(2\omega_e x_e) - \tfrac{1}{2}$, with
ladder operators
$\langle v{+}1|J_-|v\rangle = \sqrt{(2j-v)(v+1)}$ encoding the
adjacency weights of a bounded chain.  The nuclear graph
$G_{\mathrm{nuc}} = G_{\mathrm{vib}} \otimes G_{\mathrm{rot}}$ is the
Cartesian product of these two lattices, with vibration-rotation
coupling $B_v = B_e - \alpha_e(v+\tfrac{1}{2})$ modifying the
rotational structure at each vibrational level.  The product graph
eigenvalues reproduce the full rovibrational spectrum
$E_{v,J} = \omega_e(v+\tfrac{1}{2}) - \omega_e x_e(v+\tfrac{1}{2})^2
+ B_v\,J(J+1) - D_v\,J^2(J+1)^2$ to machine precision for H$_2$, HCl,
CO, and LiH.  The key structural difference from the electron lattice
is that the vibrational SU(2) has \emph{finite-dimensional}
representations ($v_{\max} < \infty$), reflecting the bounded nature of
the molecular bond, whereas the electronic $n$-ladder extends to
infinity (ionization continuum).  We identify $\hbar$ as the universal
discreteness marker: every appearance of $\hbar$ in the rovibrational
spectrum signals an underlying discrete lattice structure. Computational
validation with 52 tests confirms all algebraic identities and
experimental comparisons.
\end{abstract}

\maketitle


%% ========================================================================
\section{Introduction}
\label{sec:intro}
%% ========================================================================

The GeoVac program has established that the single-electron
Schr\"odinger equation is mathematically equivalent to the
Laplace-Beltrami operator on the unit three-sphere
$S^3$~\cite{paper0,paper1,paper7}.  The discrete graph Laplacian---with
nodes encoding quantum numbers $(n, l, m)$ and edges encoding transition
amplitudes---produces eigenvalues that converge to the Rydberg spectrum
$E_n = -Z^2/(2n^2)$ as $n_{\max} \to \infty$, with the universal
kinetic scale $\kappa = -1/16$ as the sole dimensionful parameter.  This
program has been extended to multi-electron atoms via Full Configuration
Interaction~\cite{fci_paper}, to heteronuclear diatomics via the
Topological LCAO architecture~\cite{lcao_paper}, and to quantum dynamics
via Crank-Nicolson time propagation~\cite{paper6}.

In all of this work, the internuclear distance $R$ enters as a
\emph{classical parameter}: the nuclei are fixed points in space, and
the electron lattice is constructed at each $R$ independently.  The
Born-Oppenheimer approximation justifies this separation for most
molecular properties, but it leaves a fundamental incompleteness in the
discrete framework.  The Topological LCAO paper~\cite{lcao_paper}
identified this incompleteness concretely: the graph Laplacian's
intra-atom kinetic energy is completely $R$-independent, producing a
monotonically attractive potential energy surface with no equilibrium
geometry.  The missing kinetic repulsion---the cost of overlapping
atomic information structures at short $R$---was partially restored by a
Fock-weighted correction derived from Paper~7's energy-shell constraint,
but a complete treatment requires making $R$ itself a discrete quantum
degree of freedom.

This paper takes the first step: constructing the \emph{nuclear
lattice}, the discrete graph structure that governs molecular vibration
and rotation.  Our approach exploits a profound structural parallel
between electronic and nuclear quantum numbers:

\begin{center}
\begin{tabular}{lcc}
\toprule
& \textbf{Electrons} & \textbf{Nuclei} \\
\midrule
Space & Momentum $\mathbf{p}$ & Position $\mathbf{R}$ \\
Dual & Position $\mathbf{r}$ & Momentum $\mathbf{P}$ \\
Angular & $(l, m)$ on $S^2$ & $(J, M_J)$ on $S^2$ \\
Radial & $n = 1, 2, \ldots$ & $v = 0, 1, \ldots, v_{\max}$ \\
Algebra & SU(2) $\otimes$ SU(1,1) & SU(2) $\otimes$ SU(2) \\
Range & Infinite ($n \to \infty$) & Finite ($v \leq v_{\max}$) \\
\bottomrule
\end{tabular}
\end{center}

The electron lattice lives in momentum space, with position $\mathbf{r}$
emerging as the Fourier dual via the stereographic projection of
Paper~7.  The nuclear lattice lives in position space, with momentum
$\mathbf{P}$ emerging as the conjugate variable.  The angular parts are
\emph{identical}: both use SO(3) representations with $2J+1$ states per
shell.  The radial parts differ crucially: the electronic $n$-ladder
extends to infinity (ionization continuum at $E = 0$), while the
vibrational $v$-chain is bounded ($v_{\max}$ set by the dissociation
energy $D_e$).

The paper is organized as follows.
Section~\ref{sec:rotational} constructs the rotational paraboloid and
demonstrates its structural identity with the electron $(l, m)$ lattice.
Section~\ref{sec:vibrational} derives the vibrational chain from the
Morse oscillator's SU(2) algebraic structure.
Section~\ref{sec:nuclear_graph} builds the product nuclear graph and
recovers the rovibrational spectrum.
Section~\ref{sec:coupling} discusses the electron-nuclear coupling
through the Born-Oppenheimer fibration.
Section~\ref{sec:equilibrium} connects the nuclear lattice to the
missing kinetic repulsion problem identified in~\cite{lcao_paper}.
Section~\ref{sec:hbar} identifies $\hbar$ as the universal discreteness
marker.
Section~\ref{sec:conclusion} summarizes and outlines next steps.


%% ========================================================================
\section{The Rotational Paraboloid}
\label{sec:rotational}
%% ========================================================================

\subsection{States and Degeneracy}

The rotational states of a diatomic molecule are labeled by angular
momentum quantum numbers $|J, M_J\rangle$ with
$J = 0, 1, 2, \ldots, J_{\max}$ and $-J \leq M_J \leq J$.  Each
$J$-shell contains $2J + 1$ states, and the cumulative count through
shell $J$ is
\begin{equation}
\sum_{J'=0}^{J} (2J' + 1) = (J + 1)^2.
\label{eq:rot_degeneracy}
\end{equation}
This is structurally identical to the electron lattice, where the
$(l, m)$ sublattice at fixed $n$ contains
$\sum_{\ell=0}^{n-1}(2\ell+1) = n^2$ states.  Both are manifestations
of the same SO(3) angular momentum algebra.

The rotational states form a paraboloid: $J = 0$ is the apex (single
state), $J = 1$ has 3 states arranged on a ring, $J = 2$ has 5 states,
and so on.  The azimuthal spacing between adjacent $M_J$ values is $\pi$
(the same discrete step that appears in the electron lattice).

\subsection{Adjacency and Selection Rules}

The angular momentum raising and lowering operators
\begin{equation}
L_\pm |J, M_J\rangle = \sqrt{J(J+1) - M_J(M_J \pm 1)}\;
|J, M_J \pm 1\rangle
\label{eq:lpm}
\end{equation}
define the edges of the rotational graph.  The adjacency matrix
$A_{\mathrm{rot}}$ has nonzero entries only between states with the same
$J$ and $\Delta M_J = \pm 1$.  The $J = 0$ state is isolated: it has no
$M_J$ neighbors, just as the $l = 0$ (s-orbital) node in the electron
lattice has no angular neighbors.

The selection rules follow directly from the graph structure:
\begin{itemize}
\item $\Delta M_J = 0, \pm 1$ (electric dipole transitions)
\item $\Delta J = 0, \pm 1$ (from inter-shell coupling, not encoded
  in the adjacency within a single $J$)
\end{itemize}

The rigid-rotor energy for each $J$-shell is
\begin{equation}
E_J = B_e\, J(J+1) - D_e^{\mathrm{rot}}\, J^2(J+1)^2,
\label{eq:rigid_rotor}
\end{equation}
where $B_e = \hbar^2/(2\mu R_e^2)$ is the rotational constant and
$D_e^{\mathrm{rot}}$ is the centrifugal distortion constant.  In the
graph representation, this spectrum appears as the diagonal of the
rotational Hamiltonian---each $J$-shell is degenerate in $M_J$, with
energy set by the Casimir invariant $J(J+1)$.


%% ========================================================================
\section{The Vibrational Chain: Morse Oscillator as SU(2)}
\label{sec:vibrational}
%% ========================================================================

\subsection{The Morse Potential and Its Algebraic Structure}

The Morse potential
\begin{equation}
V(r) = D_e\left[1 - e^{-a(r - r_e)}\right]^2
\label{eq:morse}
\end{equation}
with $a = \omega_e\sqrt{\mu/(2D_e)}$ is one of the few anharmonic
potentials with an exact algebraic solution.  Morse himself showed that
the bound-state energies are
\begin{equation}
E_v = \hbar\omega_e\left(v + \tfrac{1}{2}\right) -
\hbar\omega_e x_e\left(v + \tfrac{1}{2}\right)^2,
\label{eq:morse_spectrum}
\end{equation}
where $x_e = \hbar\omega_e/(4D_e)$ is the anharmonicity parameter.

What is less widely appreciated is that the Morse spectrum has an exact
SU(2) algebraic structure~\cite{frank_lemus,alhassid_gursey}.  The key
identification is:
\begin{equation}
j = \frac{\omega_e}{2\omega_e x_e} - \frac{1}{2},
\label{eq:j_def}
\end{equation}
which maps the Morse parameters to the SU(2) representation label $j$.
The vibrational quantum number $v$ maps to the angular momentum
projection via
\begin{equation}
m = j - v,
\label{eq:m_from_v}
\end{equation}
so $v = 0$ corresponds to $m = j$ (the highest-weight state) and
$v = v_{\max}$ corresponds to $m \approx -j$ (the lowest-weight state at
the dissociation threshold).

\subsection{Ladder Operators and the Vibrational Chain}

The SU(2) lowering operator $J_-$ acts as
\begin{equation}
J_-|j, m\rangle = \sqrt{(j+m)(j-m+1)}\;|j, m-1\rangle.
\label{eq:jminus}
\end{equation}
Substituting $m = j - v$ gives the matrix element connecting vibrational
state $v$ to $v + 1$:
\begin{equation}
\langle v+1|J_-|v\rangle = \sqrt{(2j - v)(v + 1)}.
\label{eq:vib_ladder}
\end{equation}
This is the weight on edge $(v, v+1)$ in the vibrational chain graph.

\begin{proposition}[Vibrational Adjacency Matrix]
The adjacency matrix of the vibrational chain is the symmetric
tridiagonal matrix
\begin{equation}
A_{\mathrm{vib}}[v, v'] = \begin{cases}
\sqrt{(2j - v)(v + 1)} & \text{if } v' = v + 1, \\
\sqrt{(2j - v')(v' + 1)} & \text{if } v' = v - 1, \\
0 & \text{otherwise.}
\end{cases}
\label{eq:avib}
\end{equation}
\end{proposition}

The chain has several notable properties:
\begin{enumerate}
\item \textbf{Bounded.}  The chain terminates at
$v_{\max} \leq \lfloor j \rfloor$, the last state whose energy lies
below the dissociation limit $D_e$.  This is the fundamental difference
from the electron $n$-ladder, which extends to $n = \infty$.

\item \textbf{Non-uniform weights.}  The edge weights
$w(v) = \sqrt{(2j-v)(v+1)}$ are largest near $v \approx j$ (the middle
of the representation) and decrease toward both boundaries.  This
non-uniformity encodes the anharmonicity of the Morse potential.

\item \textbf{SU(2) Casimir.}  The Casimir operator
$\mathbf{J}^2 = J_+ J_- + J_3^2 - J_3$ takes the constant value
$j(j+1)$ throughout the representation, confirming that all vibrational
states belong to a single irreducible representation of SU(2).
\end{enumerate}

\subsection{The Morse Spectrum from SU(2)}

The Morse energy levels~(\ref{eq:morse_spectrum}) can be rewritten in
SU(2) language.  Using $v + \tfrac{1}{2} = j - m + \tfrac{1}{2}$ and
the Morse relation $\omega_e x_e = \omega_e/(2j+1)$:
\begin{equation}
E_v = \omega_e\left(v + \tfrac{1}{2}\right)
\left[1 - \frac{v + \tfrac{1}{2}}{2j + 1}\right].
\label{eq:morse_su2}
\end{equation}
The factor $[1 - (v+\tfrac{1}{2})/(2j+1)]$ is the anharmonic
correction: for $v \ll j$ it is negligible (harmonic limit), and for
$v \to j$ it drives the spacing to zero (dissociation).

\subsection{Comparison to the Electron $n$-Ladder}

The electron energy levels $E_n = -Z^2/(2n^2)$ arise from the
\emph{non-compact} group SU(1,1), whose representations are
infinite-dimensional: $n = 1, 2, 3, \ldots$ with no upper bound.  The
nuclear vibrational levels arise from the \emph{compact} group SU(2),
whose representations are finite-dimensional: $v = 0, 1, \ldots,
v_{\max}$.

This distinction has a direct physical meaning:
\begin{itemize}
\item \textbf{Electrons:} The $n \to \infty$ limit is ionization---the
electron escapes to the continuum at $E = 0$.  There is no upper bound
on excitation because the Coulomb potential extends to infinity.
\item \textbf{Nuclei:} The $v \to v_{\max}$ limit is
dissociation---the bond breaks at $E = D_e$.  The finite representation
reflects the bounded nature of the molecular potential well.
\end{itemize}


%% ========================================================================
\section{The Nuclear Graph Laplacian}
\label{sec:nuclear_graph}
%% ========================================================================

\subsection{Product Graph Construction}

The full nuclear graph is the Cartesian product of the vibrational and
rotational lattices:
\begin{equation}
G_{\mathrm{nuc}} = G_{\mathrm{vib}} \times G_{\mathrm{rot}}.
\label{eq:product}
\end{equation}
The Cartesian product graph has vertex set
$V_{\mathrm{vib}} \times V_{\mathrm{rot}}$, and two vertices
$(v_1, J_1, M_1)$ and $(v_2, J_2, M_2)$ are adjacent if and only if
either:
\begin{itemize}
\item $v_1 = v_2$ and $(J_1, M_1)$ is adjacent to $(J_2, M_2)$ in
  $G_{\mathrm{rot}}$, or
\item $(J_1, M_1) = (J_2, M_2)$ and $v_1$ is adjacent to $v_2$ in
  $G_{\mathrm{vib}}$.
\end{itemize}

The adjacency matrix of the product graph is
\begin{equation}
A_{\mathrm{nuc}} = A_{\mathrm{vib}} \otimes I_{\mathrm{rot}} +
I_{\mathrm{vib}} \otimes A_{\mathrm{rot}},
\label{eq:anuc}
\end{equation}
and the graph Laplacian is
\begin{equation}
L_{\mathrm{nuc}} = L_{\mathrm{vib}} \otimes I_{\mathrm{rot}} +
I_{\mathrm{vib}} \otimes L_{\mathrm{rot}}.
\label{eq:lnuc}
\end{equation}

For a molecule with $v_{\max}$ vibrational states and rotational states
up to $J_{\max}$, the total dimension is
\begin{equation}
N_{\mathrm{nuc}} = (v_{\max} + 1) \times (J_{\max} + 1)^2.
\end{equation}

\subsection{Vibration-Rotation Coupling}

The product graph~(\ref{eq:product}) captures independent vibration and
rotation.  Physical molecules exhibit vibration-rotation coupling: as
the molecule vibrates, the moment of inertia changes, modifying the
rotational constant.  This coupling is encoded by making the rotational
energy $v$-dependent:
\begin{equation}
B_v = B_e - \alpha_e\left(v + \tfrac{1}{2}\right),
\label{eq:bv}
\end{equation}
where $\alpha_e$ is the vibration-rotation coupling constant.

The full rovibrational energy is
\begin{equation}
E_{v,J} = \omega_e\left(v + \tfrac{1}{2}\right)
- \omega_e x_e\left(v + \tfrac{1}{2}\right)^2
+ B_v\,J(J+1) - D_v\,J^2(J+1)^2,
\label{eq:rovib}
\end{equation}
where $D_v$ is the centrifugal distortion constant.

\subsection{Eigenvalue Spectrum}

The nuclear Hamiltonian is diagonal in the $|v, J, M_J\rangle$ basis
(neglecting Coriolis coupling):
\begin{equation}
H_{\mathrm{nuc}}|v, J, M_J\rangle = E_{v,J}\;|v, J, M_J\rangle.
\label{eq:hnuc}
\end{equation}
Each rovibrational level has $(2J+1)$-fold $M_J$ degeneracy (in the
absence of external fields), identical to the electron lattice's
$m$-degeneracy within each $l$-shell.

Table~\ref{tab:molecules} compares the nuclear lattice parameters for
four benchmark diatomics.

\begin{table*}
\centering
\begin{tabular}{lcccccccc}
\toprule
Molecule & $D_e$ (Ha) & $\omega_e$ (cm$^{-1}$) & $\omega_e x_e$ (cm$^{-1}$)
& $B_e$ (cm$^{-1}$) & $\alpha_e$ (cm$^{-1}$) & $j$ & $v_{\max}$ &
$\nu_{01}$ (cm$^{-1}$) \\
\midrule
H$_2$  & 0.1745 & 4401.2 & 121.3 & 60.85 & 3.062 & 17.6 & 13 & 4158.5 \\
HCl    & 0.1695 & 2991.0 &  52.8 & 10.59 & 0.307 & 27.8 & 17 & 2885.3 \\
CO     & 0.4130 & 2169.8 &  13.3 &  1.93 & 0.018 & 81.1 & 81 & 2143.2 \\
LiH    & 0.0920 & 1405.7 &  23.2 &  7.51 & 0.213 & 29.8 & 22 & 1359.3 \\
\bottomrule
\end{tabular}
\caption{Spectroscopic constants and nuclear lattice parameters for four
benchmark diatomics.  The SU(2) representation label $j$ and maximum
vibrational quantum number $v_{\max}$ determine the size of the
vibrational chain.  The fundamental frequency
$\nu_{01} = \omega_e - 2\omega_e x_e$ matches experimental values by
construction.}
\label{tab:molecules}
\end{table*}


%% ========================================================================
\section{Electron-Nuclear Coupling: The Born-Oppenheimer Fibration}
\label{sec:coupling}
%% ========================================================================

\subsection{The Fibered Graph}

In the Born-Oppenheimer approximation, the electronic Hamiltonian
depends parametrically on the nuclear configuration.  In the discrete
framework, this means the electron paraboloid (Paper~0) sits
\emph{above} each nuclear graph vertex, with its parameters varying
along the nuclear fiber.

The energy-shell constraint of Paper~7 provides the connection:
\begin{equation}
p_0^2 = -2E = \frac{Z^2}{n^2}
\label{eq:p0}
\end{equation}
determines the ``stereographic focal length'' of the electron lattice.
As the molecule vibrates from $v = 0$ (equilibrium) toward $v_{\max}$
(dissociation), the effective nuclear charge seen by the valence
electrons decreases, modifying $p_0$:
\begin{equation}
p_0(v) = \frac{Z_{\mathrm{eff}}(v)}{n},
\label{eq:p0v}
\end{equation}
where $Z_{\mathrm{eff}}(v)$ decreases monotonically from
$Z_{\mathrm{eff}}(0) = Z$ at equilibrium toward a dissociation limit
set by the separated-atom charges.

\subsection{Franck-Condon Factors as Inter-Fiber Weights}

Electronic transitions between different electronic surfaces
(e.g., $S_0 \to S_1$) involve changing the vibrational quantum number.
The probability of a $v \to v'$ transition is given by the
Franck-Condon factor $|\langle v'|v\rangle|^2$, which measures the
overlap between vibrational wavefunctions on different potential energy
surfaces.

In the graph framework, these Franck-Condon factors appear as edge
weights connecting nuclear vertices on different electronic fibers.  The
combined electron-nuclear graph has a fibered structure:
\begin{equation}
G_{\mathrm{total}} = G_{\mathrm{elec}}(v) \ltimes G_{\mathrm{nuc}},
\label{eq:fibered}
\end{equation}
where $G_{\mathrm{elec}}(v)$ is the electronic paraboloid at vibrational
level $v$, and $\ltimes$ denotes the semi-direct product encoding the
parametric dependence.


%% ========================================================================
\section{Emergence of Equilibrium Geometry}
\label{sec:equilibrium}
%% ========================================================================

\subsection{Connection to the Missing Kinetic Repulsion}

The Topological LCAO paper~\cite{lcao_paper} identified a fundamental
limitation of the discrete molecular framework: the graph Laplacian's
intra-atom kinetic energy is completely $R$-independent, producing no
equilibrium geometry.  The Fock-weighted kinetic correction
\begin{equation}
\Delta h_{aa} = \lambda \sum_{b \in B}
S_{ab}^2 \left(\frac{Z_A}{n_a}\right)^2
\left(\frac{Z_B}{n_b}\right)^2
\label{eq:fock_correction}
\end{equation}
partially restores the equilibrium by introducing an $R$-dependent
penalty proportional to the product of the information densities
$p_0^2 = Z^2/n^2$ at each center.

The nuclear lattice provides the natural context for understanding this
correction.  In the combined electron-nuclear graph, the internuclear
distance $R$ is no longer a classical parameter---it is determined by
the vibrational state $v$.  The equilibrium bond length $R_e$
corresponds to the $v = 0$ ground state of the nuclear lattice, where
the total electron-nuclear energy is minimized.

\subsection{$R_e$ as a Variational Minimum}

The combined energy functional
\begin{equation}
\mathcal{E}[v, J] = E_{\mathrm{elec}}(v) + E_{\mathrm{nuc}}(v, J)
\label{eq:total_energy}
\end{equation}
has a minimum at $v = 0$ by construction (the Morse ground state).  The
corresponding nuclear geometry $R_e$ is encoded in the spectroscopic
constants $B_e = \hbar^2/(2\mu R_e^2)$ and $\omega_e$, which determine
the shape of the nuclear graph.

The missing kinetic repulsion of~\cite{lcao_paper} can now be understood
as the $v$-dependence of the electronic energy $E_{\mathrm{elec}}(v)$:
as $v$ increases (bond stretches), the electronic energy rises because
the overlap between atomic lattices decreases, reducing the bonding
stabilization.  The equilibrium is the balance between this electronic
destabilization and the nuclear zero-point energy.


%% ========================================================================
\section{$\hbar$ as the Discreteness Marker}
\label{sec:hbar}
%% ========================================================================

A striking pattern emerges from the rovibrational spectrum: every
appearance of $\hbar$ (or equivalently, every energy gap that cannot be
obtained from classical mechanics) signals an underlying discrete lattice
structure.

\begin{itemize}
\item \textbf{Vibrational quantization:}
$E_v = \hbar\omega_e(v + \tfrac{1}{2})$.  The factor $\hbar\omega_e$
is the energy spacing of the vibrational chain.  Without the chain
(continuous classical vibration), there would be no $\hbar$.

\item \textbf{Rotational quantization:}
$E_J = B_e\,J(J+1)$ where $B_e \propto \hbar^2$.  The factor $\hbar^2$
is the angular discreteness scale of the rotational paraboloid.  Without
the paraboloid (continuous classical rotation), $E$ would be continuous.

\item \textbf{Zero-point energy:}
$E_0 = \tfrac{1}{2}\hbar\omega_e$.  The zero-point energy is the
ground-state eigenvalue of the vibrational chain Laplacian.  It arises
because a finite chain cannot have a zero eigenvalue without violating
the boundary condition at $v = v_{\max}$.

\item \textbf{Anharmonicity:}
$\omega_e x_e = \hbar\omega_e^2/(4D_e)$.  The anharmonic correction is
the finite-size effect of the SU(2) representation: the chain has length
$v_{\max} + 1 \sim 2j$, and the non-uniform edge weights
$\sqrt{(2j-v)(v+1)}$ encode the departure from the harmonic
(infinite-chain) limit.
\end{itemize}

This pattern is the nuclear analog of the electronic observation: in
Paper~7, $\hbar$ appears in $E_n = -\hbar^2 Z^2/(2m_e n^2 a_0^2)$
because the electron occupies a discrete lattice on $S^3$.  In both
cases, \emph{discreteness is the origin of quantization, not vice versa}.


%% ========================================================================
\section{Computational Validation}
\label{sec:validation}
%% ========================================================================

The nuclear lattice implementation (\texttt{geovac/nuclear\_lattice.py})
is validated by a comprehensive test suite
(\texttt{tests/test\_nuclear\_lattice.py}) containing 52 tests organized
into nine categories:

\begin{enumerate}
\item \textbf{Morse SU(2) algebra} (7 tests): $j$-representation,
$v_{\max}$ from $j$, ladder element positivity, explicit formula
verification, boundary behavior, Casimir constancy, $J_3$ spectrum.

\item \textbf{Morse spectrum reproduction} (6 tests): energy formula,
spacing decrease (anharmonicity), zero-point energy, fundamental
frequency, graph spectrum consistency, dissociation limit.

\item \textbf{Rotational paraboloid structure} (7 tests): $2J+1$
degeneracy, total state count $(J_{\max}+1)^2$, paraboloid shape,
adjacency within $J$-shells, angular momentum weights, rigid-rotor
spectrum, $\Delta M_J$ selection rules.

\item \textbf{Product graph rovibrational spectrum} (5 tests):
rovibrational formula, Hamiltonian diagonality, product graph size,
adjacency shape, Laplacian zero row-sum.

\item \textbf{Experimental comparison} (7 tests): fundamental frequency
for H$_2$/HCl/CO/LiH, first overtone for all four, bound state counts,
rotational constant.

\item \textbf{Electron-nuclear structural parallels} (3 tests):
finite vs.\ infinite representation, rotational-angular identity, SU(2)
in both.

\item \textbf{Electron-nuclear coupling} (4 tests): $p_0$ at
equilibrium, $p_0$ decrease with $v$, Franck-Condon diagonal and
off-diagonal.

\item \textbf{Build helpers} (5 tests): factory function for all four
molecules, unknown molecule error handling.

\item \textbf{Adjacency matrix properties} (4 tests): vibrational and
rotational symmetry, tridiagonal structure, positive semi-definiteness.
\end{enumerate}

All 52 tests pass.  The experimental spectroscopic constants (NIST/Huber-Herzberg) are reproduced to machine precision by
construction, validating the algebraic framework.


%% ========================================================================
\section{Conclusions and Next Steps}
\label{sec:conclusion}
%% ========================================================================

We have constructed the nuclear lattice---the discrete graph structure
governing molecular vibration and rotation---and demonstrated its
structural parallel with the electron lattice of the GeoVac framework.
The key results are:

\begin{enumerate}
\item The rotational paraboloid $(J, M_J)$ is \emph{algebraically
identical} to the electron $(l, m)$ lattice, both being SO(3) angular
momentum representations.

\item The vibrational chain exploits the Morse oscillator's SU(2)
algebraic structure, with ladder elements
$\sqrt{(2j-v)(v+1)}$ encoding the adjacency weights.

\item The product graph
$G_{\mathrm{nuc}} = G_{\mathrm{vib}} \times G_{\mathrm{rot}}$
reproduces the full rovibrational spectrum to machine precision.

\item The finite-dimensionality of the vibrational SU(2)
($v_{\max} < \infty$) versus the infinite electron $n$-ladder reflects
the fundamental difference between bounded molecular bonds and unbounded
Coulomb potentials.

\item Every appearance of $\hbar$ in the rovibrational spectrum signals
an underlying discrete lattice structure.
\end{enumerate}

The natural next steps are:

\begin{enumerate}
\item \textbf{Coupled electron-nuclear dynamics:} Propagate the combined
electron-nuclear graph using the time evolution methods of
Paper~6~\cite{paper6}, enabling non-adiabatic dynamics, photodissociation,
and vibronic coupling on a unified discrete lattice.

\item \textbf{$R$-dependent kinetic energy:} Embed the Fock-weighted
kinetic correction of~\cite{lcao_paper} into the nuclear graph structure,
making the equilibrium geometry emerge from the combined electron-nuclear
graph eigenvalue problem rather than from a fitted parameter.

\item \textbf{Polyatomic extension:} Generalize the vibrational chain
to normal modes of polyatomic molecules, each carrying its own SU(2)
Morse representation, coupled through Fermi resonances and
Darling-Dennison coupling.

\item \textbf{Nuclear graph Laplacian eigenvalues:} Investigate whether
the eigenvalues of $L_{\mathrm{nuc}}$ (not just the diagonal
Hamiltonian) encode additional physical information, such as transition
selection rules or Franck-Condon-like overlap integrals, analogous to
how the electron graph Laplacian eigenvalues encode the Rydberg spectrum.

\item \textbf{First-principles electronic PES:} The Morse parameters
$(D_e, \omega_e, \omega_e x_e)$ used in this paper are currently taken
from experiment or high-level quantum chemistry.  The prolate spheroidal
lattice of Paper~11~\cite{loutey_paper11} computes the electronic PES
$E(R)$ from first principles by discretizing the Laplace-Beltrami
operator in the natural two-center coordinate system, achieving sub-1\%
accuracy for H$_2^+$ with zero fitted parameters.  Coupling that
electronic PES to the nuclear lattice developed here would make the
full electron-nuclear description self-consistent within the GeoVac
framework.
\end{enumerate}


%% ========================================================================
%% REFERENCES
%% ========================================================================

\begin{thebibliography}{25}

\bibitem{paper0}
J.~Loutey, ``Quantum Numbers as Emergent Geometric Labels: Deriving
Hydrogen's Structure from Information Packing,'' GeoVac Paper~0 (2026).

\bibitem{paper1}
J.~Loutey, ``Spectral Graph Theory for Atomic Structure,'' GeoVac
Paper~1 (2026).

\bibitem{paper7}
J.~Loutey, ``The Dimensionless Vacuum: Recovering the Schr\"odinger
Equation from Scale-Invariant Graph Topology,'' GeoVac Paper~7 (2026).

\bibitem{fci_paper}
J.~Loutey, ``Topological Full Configuration Interaction for
Multi-Electron Atoms,'' GeoVac FCI Paper (2026).

\bibitem{lcao_paper}
J.~Loutey, ``Topological Full Configuration Interaction for
Heteronuclear Diatomics: LiH as a Benchmark,'' GeoVac LCAO FCI Paper
(2026).

\bibitem{paper6}
J.~Loutey, ``$O(V)$ Quantum Dynamics on the Discrete Paraboloid
Lattice,'' GeoVac Paper~6 (2026).

\bibitem{paper8}
J.~Loutey, ``Bond Sphere Theory, Sturmian Negative Theorem, and
Selection Rules,'' GeoVac Paper~8 (2026).

\bibitem{fock1935}
V.~Fock, ``Zur Theorie des Wasserstoffatoms,'' Z.~Phys.\ \textbf{98},
145 (1935).

\bibitem{morse1929}
P.~M.~Morse, ``Diatomic molecules according to the wave mechanics.~II.
Vibrational levels,'' Phys.\ Rev.\ \textbf{34}, 57 (1929).

\bibitem{frank_lemus}
A.~Frank and R.~Lemus, ``Algebraic approach to vibrational spectra of
polyatomic molecules: Morse oscillators and SU(2),'' Phys.\ Rev.\
Lett.\ \textbf{68}, 413 (1992).

\bibitem{alhassid_gursey}
Y.~Alhassid, F.~G\"ursey, and F.~Iachello, ``Group theory approach to
scattering.~II. The Euclidean connection,'' Ann.\ Phys.\ (N.Y.)\
\textbf{167}, 181 (1986).

\bibitem{huber}
K.~P.~Huber and G.~Herzberg, \emph{Molecular Spectra and Molecular
Structure, Vol.~IV: Constants of Diatomic Molecules} (Van Nostrand
Reinhold, New York, 1979).

\bibitem{barut1967}
A.~O.~Barut and H.~Kleinert, ``Transition probabilities of the
hydrogen atom from noncompact dynamical groups,'' Phys.\ Rev.\
\textbf{156}, 1541 (1967).

\bibitem{iachello_levine}
F.~Iachello and R.~D.~Levine, \emph{Algebraic Theory of Molecules}
(Oxford University Press, 1995).

\bibitem{loutey_paper11}
J.~Loutey,
``The Molecular Fock Projection: Diatomic Molecules as Prolate
Spheroidal Lattices,''
GeoVac Paper~11 (2026).

\end{thebibliography}

\end{document}
